\section{PWM DC/DC converters in DCM}

\subsection{Ctrl+F keywords: DCM, discontinuous conduction, boundary, critical \(R\), \(K\), \(\delta\), triangular inductor current, dummy load, critical load}

DCM: inductor current reaches zero before the next turn-on. The diode prevents reversal. While \(i_L=0\), the load is supplied by \(C\).

\subsection{Conduction parameter \(K\)}
\[
K=\frac{2L}{R T_s}=\frac{2L f_s}{R}.
\]
Small \(K\) (small \(L\), large \(R\), low \(f_s\)) promotes DCM.

\subsection{CCM/DCM boundary}
DCM when half the peak-to-peak ripple exceeds the CCM average inductor current, equivalently \(K<K_{\mathrm{crit}}(D)\).

\begin{center}
\begin{tabular}{@{}lll@{}}
\toprule
Topology & Boundary (DCM if) & Critical load current idea\\
\midrule
Buck & \(K<(1-D)\) & \(\Delta i_L/2>I_o\)\\
Boost & \(K<D(1-D)^2\) & \(\Delta i_L/2>I_{\mathrm{in}}\)\\
Buck-boost & \(K<(1-D)^2\) & \(\Delta i_L/2>\overline{i}_L=I_o/(1-D)\)\\
\bottomrule
\end{tabular}
\end{center}
Always-CCM for every \(D\) requires \(K\) larger than the maximum of \(K_{\mathrm{crit}}(D)\).
For boost, \(K_{\mathrm{crit}}=D(1-D)^2\) has a maximum of \(4/27\) at \(D=1/3\). Hence \(K>4/27\) keeps CCM for all \(D\).

\subsection{DCM conversion ratios (lossless, \(P_{\mathrm{in}}=P_o\))}
Inductor current is a triangle of height \(\Delta i_L\). Equate \(V_{\mathrm{in}}I_{\mathrm{in}}=V_o^2/R\).

\subsubsection{Buck DCM}
\[
\Delta i_L=\frac{V_{\mathrm{in}}-V_o}{L}D T_s,\qquad I_{\mathrm{in}}=\frac{\Delta i_L}{2}D,
\]
\[
\frac{V_o}{V_{\mathrm{in}}}=\frac{2}{1+\sqrt{1+4K/D^2}}.
\]
The extra off-conduction fraction is
\[
\delta=\frac{(V_{\mathrm{in}}-V_o)D}{V_o},\qquad D+\delta<1.
\]

\subsubsection{Boost DCM}
\[
\frac{V_o}{V_{\mathrm{in}}}=\frac{1+\sqrt{1+4D^2/K}}{2}\ge 1.
\]
\[
\delta=\frac{V_{\mathrm{in}} D}{V_o-V_{\mathrm{in}}}.
\]
At no-load, \(V_o\to\infty\); a dummy load is required. \(V_o\ge V_{\mathrm{in}}\) always.

\subsubsection{Buck-boost DCM}
\[
\frac{V_o}{V_{\mathrm{in}}}=\frac{D}{\sqrt{K}}.
\]
Approximately linear in \(D\) at fixed \(K\). Gain still spans \(0\) to \(\infty\).

\subsection{Where the CCM curve meets DCM}
On the \(M\)--\(D\) plot the straight (or \(1/(1-D)\)) CCM locus is the upper/inner boundary. The curved DCM family is parameterised by \(K\). ``Enter CCM'' marks \(K=K_{\mathrm{crit}}\).

\subsection{DCM ripple and rms (triangle of height \(\Delta i_L\))}
Capacitor current is the AC part of \(i_L\) (or of the diode current on boost/buck-boost). Peak-to-peak output ripple from the positive \(i_C\) triangle:
\[
\Delta V_o=\frac{\delta T_s(\Delta i_L-I_o)^2}{2C\,\Delta i_L}.
\]
% TODO: verify from source -- boost/buck-boost slides write the same \(\delta T_s(\Delta i_L-I_o)^2/(2C\Delta i_L)\) form.
Switch/diode rms for a triangle occupying fraction \(d_{\mathrm{cond}}\) of the period:
\[
I_{\mathrm{rms}}=\Delta i_L\sqrt{\frac{d_{\mathrm{cond}}}{3}}.
\]
Buck: transistor fraction \(D\), diode fraction \(\delta\). Average currents:
\[
I_{\mathrm{in}}=\frac{\Delta i_L D}{2},\qquad I_o=\frac{\Delta i_L(D+\delta)}{2}\quad\text{(buck; diode \(I_o=\Delta i_L\delta/2\))}.
\]
Capacitor rms:
\[
i_{C,\mathrm{rms}}^2=i_{L,\mathrm{rms}}^2-\overline{i}_L^2.
\]

Blocking voltages in DCM are the same as CCM (buck: \(V_{\mathrm{in}}\); boost: \(V_o\); buck-boost: \(V_{\mathrm{in}}+V_o\)). Some printed DCM tables list boost switch voltage as \(V_{\mathrm{in}}\); use \(V_o\).

\subsection{Exam recipe: critical \(R\) for buck DCM}
\textbf{Use when:} given \(V_s\), \(V_o\), \(L\), \(f_s\); asked for critical load resistance.

\textbf{Steps:}
\begin{enumerate}
  \item CCM duty \(D=V_o/V_s\) (valid at the CCM/DCM boundary).
  \item Boundary: \(K=(1-D)\).
  \item \(R_{\mathrm{crit}}=2L/(K T_s)=2L f_s/(1-D)\). Larger \(R\) (lighter load) \(\Rightarrow\) DCM.
\end{enumerate}

\textbf{Fast checks:} \(D<1\) so \(R_{\mathrm{crit}}>2Lf_s\).

\textbf{Common traps:} using boost's \(D(1-D)^2\); treating \(V_o/V_s\) as a DCM formula while sitting on the boundary (on the boundary CCM and DCM gains coincide).

\subsection{Exam recipe: always-CCM boost, worst-case \(D\)}
\textbf{Use when:} boost, find minimum \(I_o\) or maximum \(R\) so that CCM holds for all duty ratios.

\textbf{Steps:}
\begin{enumerate}
  \item \(K_{\mathrm{crit}}(D)=D(1-D)^2\). Maximise: \(D=1/3\), \(K_{\mathrm{max}}=4/27\).
  \item Always CCM if \(K>4/27\), i.e.\ \(R<R_{\mathrm{max}}=2L f_s/(4/27)= (27/2) L f_s\).
  \item Corresponding minimum load current uses \(V_o=V_{\mathrm{in}}/(1-D)\) at that \(D\) if a numerical \(I_o\) is required, or \(I_{o,\min}\) from \(K\) and the worst-case \(D\).
\end{enumerate}

\textbf{Fast checks:} \(4/27\approx0.148\).

\textbf{Common traps:} maximising \((1-D)^2\) instead of \(D(1-D)^2\); mixing \(I_{\mathrm{in}}\) and \(I_o\).

\subsection{Exam recipe: design \(L\) and \(C\) of CCM buck-boost}
\textbf{Use when:} specs \(V_{\mathrm{in}}\), \(V_o\), \(I_o\), \(\Delta i_L\), \(\Delta V_o\), \(f\).

\textbf{Steps:}
\begin{enumerate}
  \item \(D=V_o/(V_{\mathrm{in}}+V_o)\) (magnitudes; inverted output).
  \item \(L=V_{\mathrm{in}} D T_s/\Delta i_L\).
  \item \(C=I_o D T_s/\Delta V_o\) because \(\Delta V_o=I_o D T_s/C\).
\end{enumerate}

\textbf{Fast checks:} \(V_{\mathrm{in}}=5\,\mathrm{V}\), \(V_o=12\,\mathrm{V}\) \(\Rightarrow\) \(D=12/17\).

\textbf{Common traps:} using buck \(D=V_o/V_{\mathrm{in}}\); using \(8LC\) buck ripple on a buck-boost.

\subsection{Exam recipe: buck-boost DCM onset and DCM \(\Delta V_o\)}
\textbf{Use when:} given CCM-looking numbers, asked how far \(R\) must move for DCM, then DCM ripple.

\textbf{Steps:}
\begin{enumerate}
  \item At the CCM gain, \(D=V_o/(V_{\mathrm{in}}+V_o)\).
  \item \(K_{\mathrm{crit}}=(1-D)^2\), \(R_{\mathrm{crit}}=2L f_s/K_{\mathrm{crit}}\). DCM for \(R>R_{\mathrm{crit}}\) (lighter load). The wording ``how much is the load resistance dropped'' in slides is inconsistent with DCM at light load; use \(R>R_{\mathrm{crit}}\).
  \item In DCM, \(\Delta i_L=V_{\mathrm{in}} D T_s/L\), \(\delta=D V_{\mathrm{in}}/V_o\), then \(\Delta V_o=\delta T_s(\Delta i_L-I_o)^2/(2C\Delta i_L)\).
\end{enumerate}

\textbf{Common traps:} applying CCM \(\Delta V_o=D V_o T_s/(RC)\) after the converter has entered DCM.

\subsection{Recognition rules}
Three inductor-voltage levels (on, off-conducting, idle) \(\Rightarrow\) DCM. Two levels \(\Rightarrow\) CCM. Gain depends on \(R\) in DCM (buck, boost) except the compact buck-boost form \(D/\sqrt{K}\) which still contains \(R\) through \(K\).

\subsection{One-line checks}
On the CCM/DCM boundary the two \(M\) expressions coincide. At fixed \(D\), DCM buck \(M<D\) and DCM boost \(M>1/(1-D)\). Boost DCM at \(I_o\to 0\) makes \(V_o\) run away.
